\documentclass[11pt]{article}
\usepackage[utf8]{inputenc}
\usepackage{amsmath,amssymb,amsthm}
\usepackage[margin=1in]{geometry}
\usepackage{algorithm,algpseudocode}

\newtheorem{theorem}{Theorem}
\newtheorem{lemma}[theorem]{Lemma}
\newtheorem{corollary}[theorem]{Corollary}
\newtheorem{definition}[theorem]{Definition}

\title{Problem 173: Using Up to One Million Tiles How Many Different\\``Hollow'' Square Laminae Can Be Formed?}
\author{Project Euler Solution}
\date{}

\begin{document}
\maketitle

\section{Problem Statement}

A hollow square lamina is formed by removing a smaller concentric square hole from a larger square. Using up to $N = 10^6$ unit-square tiles, how many distinct hollow square laminae can be formed?

\section{Mathematical Development}

\begin{definition}
A \emph{hollow square lamina} is parameterized by $(a,b)$ with $a > b \geq 1$, $a \equiv b \pmod{2}$, where $a$ is the outer side length and $b$ the inner side length.
\end{definition}

\begin{theorem}[Tile count]
A hollow square lamina $(a,b)$ uses exactly $t(a,b) = a^2 - b^2$ tiles.
\end{theorem}

\begin{proof}
The lamina consists of $a^2$ cells minus the $b^2$ removed interior cells.
\end{proof}

\begin{lemma}[Factored form]
With border width $k = (a{-}b)/2 \geq 1$: $t = 4k(a{-}k)$. For fixed $k$, $a$ ranges over $a \geq 2k{+}1$ subject to $4k(a{-}k) \leq N$.
\end{lemma}

\begin{proof}
$t = (a{-}b)(a{+}b) = 2k(2a{-}2k) = 4k(a{-}k)$. Condition $b \geq 1$ gives $a \geq 2k+1$.
\end{proof}

\begin{lemma}[Range of $a$]
The minimum tile count for outer side $a$ is $t_{\min}(a) = 4(a{-}1)$ (at $b = a{-}2$). Hence $a \leq N/4 + 1$.
\end{lemma}

\begin{proof}
$a^2 - (a{-}2)^2 = 4a - 4 = 4(a{-}1)$. Requiring $4(a{-}1) \leq N$ gives $a \leq N/4 + 1$.
\end{proof}

\begin{theorem}[Counting formula]
For each $a \geq 3$ with $4(a{-}1) \leq N$, define $b_{\max} = a - 2$ and $b_{\min}$ as the smallest positive integer with $b_{\min} \equiv a \pmod{2}$ and $b_{\min}^2 \geq a^2 - N$. The number of valid laminae with outer side $a$ is
\[
\operatorname{count}(a) = \left\lfloor\frac{b_{\max} - b_{\min}}{2}\right\rfloor + 1
\]
provided $b_{\min} \leq b_{\max}$, and zero otherwise.
\end{theorem}

\begin{proof}
Valid values of $b$ form the arithmetic progression $\{b_{\min}, b_{\min}{+}2, \ldots, b_{\max}\}$ with common difference~2. Its cardinality is $\lfloor(b_{\max} - b_{\min})/2\rfloor + 1$.
\end{proof}

\begin{corollary}
The total is $\sum_{a=3}^{\lfloor N/4+1\rfloor}\operatorname{count}(a)$.
\end{corollary}

\section{Algorithm}

\begin{algorithmic}[1]
\State $N \gets 10^6$;\; $\text{total} \gets 0$
\For{$a = 3, 4, \ldots$ while $4(a{-}1) \leq N$}
    \State $b_{\max} \gets a - 2$
    \State $b_{\min} \gets \lceil\sqrt{\max(1,\, a^2 - N)}\,\rceil$, adjusted to satisfy $b_{\min} \equiv a \pmod{2}$
    \If{$b_{\min} \leq b_{\max}$}
        \State $\text{total} \mathrel{+}= \lfloor(b_{\max} - b_{\min})/2\rfloor + 1$
    \EndIf
\EndFor
\State \Return total
\end{algorithmic}

\section{Complexity Analysis}

\begin{itemize}
    \item \textbf{Time:} $O(N/4)$ iterations, each $O(1)$. Total: $O(N)$.
    \item \textbf{Space:} $O(1)$.
\end{itemize}

\section{Answer}

\[
\boxed{1572729}
\]

\end{document}
